\textbf{Agent capability and security benchmarks.} A capability-side line asks whether agents can complete multi-step tasks: AgentBench \citep{liu2024agentbench} spans eight environments, and $\tau$-bench \citep{yao2024taubench} evaluates whether an agent respects written domain policies---the closest analogue to our setting, but with a cooperative user, where policy adherence is a component of \emph{utility} rather than a security property. The security-side literature established that tool-using agents inherit an attack surface from the data they read \citep{greshake2023injection, perez2022ignore}: ToolEmu \citep{ruan2024toolemu} emulates and LM-grades tool risk; InjecAgent \citep{zhan2024injecagent} enumerates indirect-injection cases; AgentHarm \citep{andriushchenko2025agentharm} measures jailbroken execution of malicious tasks; R-Judge \citep{yuan2024rjudge} tests after-the-fact risk recognition; and AgentDojo \citep{debenedetti2024agentdojo}, our closest methodological relative, scores utility and security with programmatic checks. We differ in two respects (Table~\ref{tab:related}). First, in what counts as a successful attack: our attacks---a structured payment, a routed approval---are each individually well-formed and produce a transcript that reads like competent task completion, detectable only as a predicate over final state, which is why grading is on deterministic predicates and not output judgments. Second, in what the design attributes: we hold the scenario fixed and vary only the control posture (\textit{none}/\textit{advisory}/\textit{enforced}), turning a single ASR into a decomposition of where safety came from.

\begin{table}[h]
\centering
\small
\caption{Positioning against related agent benchmarks. ``Posture axis'' means the same scenario is replayed under varying enforcement strength so that model judgment and system control can be attributed separately. ``Regulatory trace'' means each threat is mapped to a named obligation and the control that discharges it.}
\label{tab:related}
\begin{tabular}{|p{2.3cm}|p{2.5cm}|p{2.9cm}|p{1.9cm}|p{1.6cm}|p{1.5cm}|}
\hline
\textbf{Benchmark} & \textbf{Domain} & \textbf{Attack delivery} & \textbf{Environment} & \textbf{Posture axis} & \textbf{Reg.\ trace} \\
\hline
AgentBench & General (8 environments) & None (capability only) & Executed & No & No \\
\hline
$\tau$-bench & Retail, airline & None (cooperative user; policy in prompt) & Executed, mutable & No & No \\
\hline
ToolEmu & General tool use & Risky instructions & LM-emulated & No & No \\
\hline
AgentDojo & Productivity suites (incl.\ banking) & Indirect injection & Executed, mutable & No & No \\
\hline
InjecAgent & Tool-integrated agents & Indirect injection & Tool-call level & No & No \\
\hline
AgentHarm & General harm (11 categories) & Direct malicious request $+$ jailbreak & Synthetic tools & No & No \\
\hline
FinVault & Financial & Injection, jailbreak, impersonation & Executed, state-writable & No & Case-driven \\
\hline
\textbf{This work} & Financial & Indirect injection only & Executed, state-writable & \textbf{Yes (3)} & \textbf{Predicate-level} \\
\hline
\end{tabular}
\end{table}

Two entries deserve a note: AgentDojo includes a banking suite but ties no task to a regulatory obligation and varies no enforcement, and FinVault is grounded in regulatory cases at the scenario level (hence ``case-driven''), whereas our mapping runs all the way to the scored predicate (\S\ref{sec:regmap}). That decomposition connects our results to the defense literature: model-level defenses train a priority ordering over instruction sources \citep{wallace2024hierarchy}, while system-level defenses argue injection should be defeated by construction, outside the model \citep{debenedetti2025camel}. Our advisory-versus-enforced gap directly measures the premise behind the latter---stated policy is worth $+71$ points for GPT-4o but only $+55$ for GPT-4o-mini (which still ends at 28\%), while enforcement is worth whatever remains for every model without exception. The variance in the advisory column is a property of the model; the invariance of the enforced column is a property of the system.

\textbf{Financial LLM evaluation.} Prior financial-LM work concentrated on capability \citep{wu2023bloomberggpt} and, recently, content-level safety: FinSafetyBench \citep{hou2026finsafetybench} measures whether a model will \emph{say} something it should refuse---a framing that does not reach our failure mode, where the model says nothing objectionable and simply calls two tools in a control-defeating order. Concurrent work FinVault \citep{yang2026finvault} evaluates financial agent safety in execution-grounded sandboxes, reporting ASR up to 50\% under injection, jailbreaking, and impersonation; its headline (defenses transfer poorly to financial agents) is consistent with our 71\% no-policy GPT-4o result. The two are complementary: FinVault covers breadth (31 scenarios, 107 vulnerabilities, 963 test cases) and answers ``which attacks land,'' while we cover the control dimension (48 scenarios replayed under three postures, three trials, every rate a 95\% Wilson interval \citep{wilson1927probable} with permutation-tested differences \citep{ernst2004permutation}) and answer ``which layer of the stack stopped them.'' The controls we test are not invented for the benchmark: reporting thresholds and anti-structuring are statutory \citep{usc5324, cfr1010311}, and segregation of duties and dual authorization are standard internal-control practice \citep{sox2002}---specified in the audit literature as who may authorize what under which conditions, precisely the form a deterministic state predicate takes.
